\documentclass{jsarticle}
\usepackage{amsmath}
\usepackage{amssymb}
\begin{document}
\title{Jones多項式からの各株式発行からの需要と供給、\\
株価の分配関数、多項式からナスダックへの金融機構の \\
株式発行のマクロとミクロの波及と流れ}
\author{高島彩さんの成蹊大学での経済理論の研究論文}
\date{}
\maketitle

   Jones多項式から各会社の株価への値打ちから日経平均株価への
   株価の価値の波及、東証株価指数Topixの全銘柄からの全体の会社の
   指数からの平均株価の決定量、そして、
   アメリカの株式市場でのダウ平均株価の終値の収束
   (戦略、reco level theory,株式会社への)全体（株式会社
   への）需要と供給。始値から終値への分配、このJones多項式の出力から
   終値への各同値類と剰余類からの部分群論の、各平均株価の推移率の予測、
   この応用は、１９２９年の世界恐慌がニューヨークダウ平均株価の暴落が、
   このJones多項式からの流れが、終値で破綻して、金融の共鳴力が、
   干渉してしまい、カタストロフィー現象として、金融の流れが乱れた結果で
   あると、この論文から分析できる。
   $$ \Hat{V} = {1 \over \sqrt{t}} \vec{V} - {1 \over t}\Hat{V} $$
   $$ \Hat{V} = \sqrt{t}\Hat{V} - t \Hat{V} $$
   $$ \Hat{V}_k = \sum_{\vec{\sigma}} \langle K| \vec{\sigma}\rangle
   \left(\sqrt{t} + {1 \over \sqrt{t}} \right)^{||\vec{\sigma}||} $$
   $$ {1 \over \sqrt{t}}\vec{V}, -{1 \over t}\Hat{V},\\ 
   \sqrt{t}\Hat{V}, -t \vec{V} $$
   この４種類の式から、reco level theory
   が、アダム・スミスの神の見えざる手を経済理論のミクロとマクロ経済の
   具現化を実現させている。
   物理では、$$ {d \over df}F = m(x) = e^f + e^{-f} = 2 i \sin (i x \log x) $$
   と、空間の仕組みが経済理論の場の理論にもなっている。周期関数と
   ピンポイントリサーチが合わさっている。
   これらの式から、前半と後半の株式市場の流れが、シンメトリーとなって
   いるために、後半の$ e^{-f} $が破綻すると、シンメトリーのために、
   前半の$ e^f $も破綻している状態になっている。全体が機能しないと、
   金融の流れが乱れていく。この破綻を回避するには、破綻が連鎖されている
   金融の周期をこれらの式から、調節が可能であることも、類推できる。
   破綻した周期を、正常であった周期に重ね合わせの原理で、ピンポイントリサーチ
   で待って、金融の投資を正常な周期関数で合わせて、元の正常な周期に
   これらの式から戻せることが可能と推測できる。
\end{document}

